% paper4_grand_unification_outline.tex
% One Equation, Three Scales: Sigma = D(rho_spacetime || rho_matter)
% as the Universal Action
% Author: Sheng-Kai Huang
% Status: CONCEPTUAL OUTLINE
% Compile with: pdflatex paper4_grand_unification_outline.tex

\documentclass[prl,twocolumn,superscriptaddress,longbibliography,floatfix]{revtex4-2}

\usepackage{amsmath}
\usepackage{amssymb}
\usepackage{amsfonts}
\usepackage{amsthm}
\usepackage{bm}
\usepackage{hyperref}
\usepackage{booktabs}
\usepackage{graphicx}
\usepackage{xcolor}

\newtheorem{theorem}{Theorem}

% Convenience macros (same as Paper 1)
\newcommand{\kB}{k_{\mathrm{B}}}
\newcommand{\Tr}{\mathrm{Tr}}
\newcommand{\id}{\mathrm{id}}
\newcommand{\Petz}{\mathcal{R}}
\newcommand{\chan}{\mathcal{N}}
\newcommand{\hilb}{\mathcal{H}}
\newcommand{\code}{\mathcal{C}}
\newcommand{\KL}{D}
\newcommand{\ket}[1]{|#1\rangle}
\newcommand{\bra}[1]{\langle#1|}
\newcommand{\braket}[2]{\langle#1|#2\rangle}
\newcommand{\op}[2]{|#1\rangle\langle#2|}
\newcommand{\abs}[1]{\left|#1\right|}
\newcommand{\norm}[1]{\left\|#1\right\|}
\newcommand{\tauasym}{\tau}

% Outline-specific formatting
\newcommand{\known}{\textsc{[known]}}
\newcommand{\newsyn}{\textsc{[new synthesis]}}
\newcommand{\newform}{\textsc{[new formula]}}
\newcommand{\spec}{\textsc{[speculative]}}
\newcommand{\towrite}[1]{{\color{blue}\textbf{[TO WRITE:} #1\textbf{]}}}
\newcommand{\toderive}[1]{{\color{red}\textbf{[TO DERIVE:} #1\textbf{]}}}
\newcommand{\assumed}[1]{{\color{orange}\textbf{[ASSUMED:} #1\textbf{]}}}

\begin{document}

\title{One Equation, Three Scales:\\
$\Sigma = D(\rho_{\mathrm{spacetime}} \| \rho_{\mathrm{matter}})$
as the Universal Action}

\author{Sheng-Kai Huang}
\email{akai@fawstudio.com}
\affiliation{Independent Researcher}

\date{\today}

\begin{abstract}
\towrite{Final abstract after Papers 1--3 are complete.}

Papers~1--3 of this series demonstrated that the entropy production
$\Sigma = D(\rho\|\sigma) - D(\chan(\rho)\|\chan(\sigma))$ governs
three seemingly unrelated domains: quantum information recovery
(Paper~1), strong-field gravity near black holes (Paper~2), and
galactic rotation curves (Paper~3). Here we propose that these are
not three separate applications but three solutions of
\emph{one equation}:
\begin{equation*}
\Sigma = D(\rho_{\mathrm{spacetime}} \| \rho_{\mathrm{matter}}),
\end{equation*}
where $\rho_{\mathrm{spacetime}}$ and $\rho_{\mathrm{matter}}$ are
density matrices encoding the geometry and matter content,
respectively. We argue that three mathematical
realizations---modular flow, gravitational Landauer principle, and
quantum channel analysis---all converge on the universal identity
$\Sigma_{\mathrm{grav}} = -\ln(-g_{00})$, independent of any
specific action principle. Bianconi's gravity-from-entropy framework
(PRD~2025) provides one formal setting for
$\Sigma = D(\rho_{\mathrm{spacetime}} \| \rho_{\mathrm{matter}})$,
though its vacuum solutions give Schwarzschild $+ O(r^{-4})$
corrections~\cite{Thattarampilly2026}, not the exponential metric.
The exponential metric instead follows from promoting the identity
to an exact equation. Different boundary conditions yield different
solutions: strong-field ($r \sim r_s$) $\to$ exponential metric,
weak-field ($r \sim \mathrm{kpc}$) $\to$ running $G$ and flat
rotation curves, no-gravity limit $\to$ standard quantum
information theory. We identify the matching conditions between
regimes, derive the emergent arrow of time
$\tauasym = 1 - \exp(-\Sigma/2)$ as a universal consequence,
and discuss open problems including the Born rule, Planck-scale
physics, and cosmological initial conditions.
\end{abstract}

\maketitle

%=======================================================================
\section{The Three Scales}
\label{sec:three_scales}
%=======================================================================

\subsection*{Content outline}

\newsyn\ \emph{Summary of Papers~1--3.}---The first three papers
of this series established that the entropy production
$\Sigma$ governs physics at three different scales:

\begin{table}[htbp]
\caption{\label{tab:three_scales}The three scales unified by
$\Sigma$.}
\begin{ruledtabular}
\begin{tabular}{lccc}
& \textbf{Paper 1} & \textbf{Paper 2} & \textbf{Paper 3} \\
\colrule
Domain & QI & Strong field & Weak field \\
Scale & any & $r \sim r_s$ & $r \sim \mathrm{kpc}$ \\
$\Sigma$ & generic & $r_s/r$ & $r_s/r + \alpha\ln r$ \\
$\chan$ & any CPTP & grav.\ redshift & RG-improved grav. \\
Key result & $F \geq e^{-\Sigma/2}$ & exp.\ metric & flat curves \\
Status & \known & \newsyn & \newsyn \\
\end{tabular}
\end{ruledtabular}
\end{table}

\newsyn\ \emph{The pattern.}---Each paper uses the \emph{same}
mathematical structure:
\begin{enumerate}
\item A quantum channel $\chan$ that processes information;
\item The QRE drop $\Sigma$ under $\chan$;
\item The Petz recovery map as the optimal retrodiction;
\item The fidelity bound $F \geq \exp(-\Sigma/2)$.
\end{enumerate}
The only difference is the \emph{channel}: a general CPTP map
(Paper~1), the gravitational redshift channel (Paper~2), or the
RG-improved gravitational channel (Paper~3). This suggests that
all three are limits of a single, more fundamental equation.

\emph{The $F = ma$ analogy.}---Newton's second law $F = ma$ is one
equation. With different force laws and boundary conditions:
\begin{itemize}
\item Gravity + central field $\to$ Kepler orbits (planets);
\item Hooke's law + spring $\to$ simple harmonic oscillation;
\item Viscosity + continuum $\to$ Navier--Stokes (fluids).
\end{itemize}
Nobody says there are ``three different $F = ma$'s.'' There is
one equation with three classes of solutions. We propose the same
for $\Sigma$.

%=======================================================================
\section{The Unified Equation}
\label{sec:unified}
%=======================================================================

\subsection*{Content outline}

\spec\ \emph{The proposal.}---We conjecture that gravity, quantum
information, and the arrow of time are all governed by
\begin{equation}
\boxed{
\Sigma = D(\rho_{\mathrm{spacetime}} \| \rho_{\mathrm{matter}}).
}
\label{eq:unified}
\end{equation}

\emph{Interpretation.}---$\rho_{\mathrm{spacetime}}$ encodes the
geometry (metric, curvature) as a density matrix.
$\rho_{\mathrm{matter}}$ encodes the matter content (stress-energy)
as a density matrix. The quantum relative entropy between them is
the entropy production $\Sigma$: the ``distance'' between what
spacetime ``is'' and what matter ``requires'' it to be.

\known\ \emph{Formal framework: Bianconi (2025).}---Bianconi~\cite{Bianconi2025}
provided the mathematical framework in which
Eq.~\eqref{eq:unified} is well defined. In the gravity-from-entropy
(GfE) framework:
\begin{itemize}
\item The metric $g_{\mu\nu}$ is promoted to a density matrix
  $\rho_g$ on the space of field configurations;
\item The matter-induced metric (the metric that matter ``wants''
  via stress-energy) is another density matrix $\rho_m$;
\item The gravitational action is the QRE:
  $S_{\mathrm{grav}} = D(\rho_g \| \rho_m)$;
\item The Einstein equations arise from extremizing $D(\rho_g \| \rho_m)$;
\item In the low-coupling limit, this reduces to the standard
  Einstein--Hilbert action.
\end{itemize}

\known\ \emph{Bianconi's key results}~\cite{Bianconi2025,BianconiCosmology2025}:
\begin{itemize}
\item Modified Einstein equations that reduce to standard GR at
  low coupling;
\item Predicts small positive cosmological constant;
\item The ``$G$-field'' (extra degree of freedom from QRE structure)
  is a candidate for dark matter;
\item FRW cosmological solutions exist and are thermal~\cite{BianconiCosmology2025}.
\end{itemize}

\newsyn\ \emph{Connection to our framework.}---Bianconi's $S_{\mathrm{grav}}
= D(\rho_g \| \rho_m)$ is \emph{literally} our
$\Sigma = D(\rho_{\mathrm{spacetime}} \| \rho_{\mathrm{matter}})$.
The GfE framework provides the formal definition; our series
provides the physical interpretation: $\Sigma$ is the entropy
production, $\tauasym = 1 - \exp(-\Sigma/2)$ is the arrow of time,
and the Petz recovery map is the optimal retrodiction.

\spec\ \emph{Critical caveat: GfE vacuum solutions.}---Thattarampilly
et~al.~\cite{Thattarampilly2026} solved the full GfE field equations
for a static, spherically symmetric vacuum and found
$g_{00} = -(1 - r_s/r - \beta r_s^2/(48r^4))$: Schwarzschild with
$O(r^{-4})$ perturbative corrections, \emph{not} the exponential
metric (which has $O(r^{-2})$ corrections). This means GfE in its
current form does not produce the exponential metric as a vacuum
solution. Our unified equation $\Sigma = D(\rho_{\mathrm{spacetime}}
\| \rho_{\mathrm{matter}})$ should therefore be understood as
a \emph{principle}---the entropy production equals the QRE between
geometry and matter---rather than being identified with any
specific action functional.

\newsyn\ \emph{Three routes independent of GfE.}---The identity
$\Sigma_{\mathrm{grav}} = -\ln(-g_{00})$ is supported by three independent
derivations that do not depend on any specific action principle:
\begin{enumerate}
\item \emph{Modular flow}: The fractional QRE loss computed via
  Dorau--Much modular theory~\cite{DorauMuch2025} gives
  $(D_{\mathrm{in}} - D_{\mathrm{out}})/D_{\mathrm{in}} = r_s/r$.
\item \emph{Gravitational Landauer}: The Tolman-modified erasure
  cost~\cite{Herrera2020} gives $\Sigma = -\ln(-g_{00})$.
\item \emph{Quantum channel}: A pure-loss bosonic channel with
  $\eta = -g_{00}$ gives $\Sigma = -\ln\eta$; Petz recovery for
  this channel has been experimentally
  demonstrated~\cite{ChenSongScarani2025}.
\end{enumerate}
The exponential metric follows from promoting the identity
$\sqrt{-g_{00}} = e^{-\Sigma/2}$ to an exact equation, not from
solving modified field equations.

%=======================================================================
\section{Strong-Field Solution}
\label{sec:strong_field}
%=======================================================================

\subsection*{Content outline}

\newsyn\ \emph{Boundary conditions.}---Static, spherically symmetric,
$r \sim r_s$, asymptotically flat.

\emph{From Paper~2.}---Under these conditions, the entropy production
is~\cite{Paper2}
\begin{equation}
\Sigma_{\mathrm{strong}} = \frac{r_s}{r} = \frac{2GM}{c^2 r},
\label{eq:sigma_strong}
\end{equation}
and the Petz recovery fidelity gives the metric component
\begin{equation}
g_{00} = -F^2 = -\exp\!\left(-\frac{r_s}{r}\right).
\label{eq:exp_metric}
\end{equation}
This is the ``exponential metric'' that replaces Schwarzschild:
\begin{itemize}
\item No event horizon (the exponential never reaches zero);
\item All curvature invariants remain finite;
\item Testable predictions: GW echoes, neutron star redshift,
  QNM frequencies, shadow size, ISCO radius.
\end{itemize}

\toderive{Show that Eq.~\eqref{eq:sigma_strong} follows from
Eq.~\eqref{eq:unified} under the stated boundary conditions.
Specifically: solve $\delta D(\rho_g \| \rho_m) = 0$ for a
static, spherically symmetric vacuum with mass $M$ at the origin.
The solution should be $g_{00} = -\exp(-r_s/r)$, reproducing
Paper~2.}

\toderive{Resolve the tension: GfE gives Schwarzschild $+ O(r^{-4})$
  corrections~\cite{Thattarampilly2026}, while our identity gives the
  exponential metric. Three possible resolutions:
  (a)~The exponential metric is the correct vacuum geometry and GfE
  needs modification;
  (b)~The true vacuum geometry is Schwarzschild $+$ corrections, and the
  exponential metric is an idealization;
  (c)~The unified equation $\Sigma = D(\rho_{\mathrm{st}} \| \rho_{\mathrm{m}})$
  requires a different mathematical realization than GfE.
  This is the central open problem of the program.}

\emph{Matching condition at large $r$.}---As $r \to \infty$,
$\Sigma_{\mathrm{strong}} \to 0$ and $g_{00} \to -1$: flat
spacetime. The strong-field solution must match onto the weak-field
solution (Paper~3) at intermediate radii $r \gg r_s$ but
$r \ll r_c \approx 37$~kpc.

%=======================================================================
\section{Weak-Field Solution}
\label{sec:weak_field}
%=======================================================================

\subsection*{Content outline}

\newsyn\ \emph{Boundary conditions.}---Quasi-static, approximately
axisymmetric (disk galaxy), $r \sim \mathrm{kpc}$, weak-field
($\Phi/c^2 \ll 1$).

\emph{From Paper~3.}---Under these conditions, the RG flow of the
gravitational channel introduces quantum corrections~\cite{Paper3}:
\begin{equation}
\Sigma_{\mathrm{weak}}(r) = \frac{r_s}{r}
  + \frac{4k_* r_s}{\pi}\,\ln\!\left(\frac{r}{r_0}\right),
\label{eq:sigma_weak}
\end{equation}
where the logarithmic term arises from the marginal running
($\eta = 1$) of Newton's constant $G(k)$ at IR scales
$k \ll k_*$.

\emph{Physical consequence.}---The logarithmic correction produces
an additional $1/r$ force that dominates at $r \gg r_c = 1/k_*$,
giving flat rotation curves without dark matter particles.

\toderive{Show that Eq.~\eqref{eq:sigma_weak} follows from
Eq.~\eqref{eq:unified} under weak-field boundary conditions with
quantum corrections included. This requires:
(i)~solving $\delta D(\rho_g \| \rho_m) = 0$ at the semiclassical
level (reproducing Dorau--Much);
(ii)~including one-loop corrections via the DPI/RG connection
(Casini--Huerta);
(iii)~showing that the marginal running $\eta = 1$ emerges
naturally.}

\emph{Matching to strong field.}---At $r \gg r_s$ but $r \ll r_c$,
the logarithmic correction is negligible and
$\Sigma_{\mathrm{weak}} \approx r_s/r$, matching
$\Sigma_{\mathrm{strong}}$ smoothly. The two solutions are not
disjoint---they are limits of the same equation.

%=======================================================================
\section{Quantum Information Solution}
\label{sec:qi}
%=======================================================================

\subsection*{Content outline}

\known\ \emph{Boundary conditions.}---No gravity ($G = 0$), general
CPTP map $\chan$, finite-dimensional Hilbert spaces.

\emph{From Paper~1.}---With no gravity, $\Sigma$ reduces to the
generic entropy production of any quantum channel~\cite{Paper1}:
\begin{equation}
\Sigma_{\mathrm{QI}} = D(\rho\|\sigma) - D(\chan(\rho)\|\chan(\sigma)),
\label{eq:sigma_qi}
\end{equation}
and the Petz recovery map $\Petz_{\sigma,\chan}$ provides the
optimal retrodiction with fidelity bound
$F \geq \exp(-\Sigma/2)$.

\newsyn\ \emph{The gravity $\to$ QI limit.}---When gravity is
turned off, $\rho_{\mathrm{spacetime}} \to \rho_{\mathrm{flat}}$
(Minkowski geometry) and the QRE
$D(\rho_{\mathrm{spacetime}} \| \rho_{\mathrm{matter}})$ reduces to
the matter-sector entropy production alone. The gravitational
contribution vanishes: $\Sigma_{\mathrm{grav}} = 0$,
$\tauasym_{\mathrm{grav}} = 0$. Only the channel-induced entropy
production remains.

\emph{Connection to quantum error correction.}---From Paper~1,
$\Sigma = 0$ $\Leftrightarrow$ $\tauasym = 0$ $\Leftrightarrow$
perfect Petz recovery $\Leftrightarrow$ quantum Markov chain
$\Leftrightarrow$ Knill--Laflamme conditions. QEC is the
deliberate construction of a subsystem where $\Sigma = 0$:
an artificial elimination of the arrow of time within a local
code space.

%=======================================================================
\section{Matching Conditions and Predictions}
\label{sec:matching}
%=======================================================================

\subsection*{Content outline}

\newsyn\ \emph{The three regimes are connected.}---The solutions of
Eq.~\eqref{eq:unified} in different boundary conditions must
\emph{match} at their boundaries:

\begin{enumerate}
\item \emph{Strong $\to$ weak}: At $r \gg r_s$, the exponential
  metric linearizes to Newton + small corrections. The logarithmic
  correction from Paper~3 emerges as the leading quantum correction
  to this classical limit.

\item \emph{Weak $\to$ QI}: At scales where $\Phi/c^2 \to 0$
  (laboratory), the gravitational $\Sigma$ becomes negligible and
  only quantum-channel entropy production remains.

\item \emph{Weak $\to$ cosmological}: At $r \sim r_H = c/H_0$
  (Hubble radius), the de~Sitter background contribution
  $\Sigma_{\mathrm{dS}} \sim r/r_H$ becomes dominant, connecting
  galactic dynamics to cosmic expansion.
\end{enumerate}

\spec\ \emph{Intermediate-scale predictions.}---At scales where
two regimes overlap, \emph{both} corrections matter simultaneously.
Specific predictions:
\begin{itemize}
\item At $r \sim 100$~kpc (galaxy cluster outskirts): both
  running-$G$ and de~Sitter corrections contribute. This may
  resolve the Bullet Cluster puzzle if the two effects combine
  to produce the observed mass-light separation.
\item At $r \sim r_s \cdot k_*/k_s$ (strong-field with quantum
  corrections): the exponential metric acquires logarithmic
  modifications. Testable via precision measurements of
  neutron star properties.
\end{itemize}

\spec\ \emph{The three-term approximation.}---In the appropriate
limits, the unified $\Sigma$ can be approximated as
\begin{equation}
\Sigma(r) \approx \underbrace{\frac{r_s}{r}}_{\text{local gravity}}
  + \underbrace{\alpha\,\ln\!\left(\frac{r}{r_c}\right)}_{\text{RG running}}
  + \underbrace{\beta\,\frac{r}{L_{\mathrm{dS}}}}_{\text{de Sitter}},
\label{eq:three_term}
\end{equation}
where:
\begin{itemize}
\item Term 1 dominates near compact objects ($r \lesssim r_s$);
\item Term 2 dominates at galactic scales ($r_c \lesssim r \lesssim L_{\mathrm{dS}}$);
\item Term 3 dominates at cosmological scales ($r \sim L_{\mathrm{dS}}$).
\end{itemize}
\textbf{This is an approximation}, not the fundamental equation.
The fundamental equation is Eq.~\eqref{eq:unified}; the three
terms emerge as limits.

%=======================================================================
\section{The Arrow of Time as Emergent}
\label{sec:time_arrow}
%=======================================================================

\subsection*{Content outline}

\newsyn\ \emph{The central thesis.}---The arrow of time is not a
separate law of nature. It \emph{is} the entropy production $\Sigma$:
\begin{equation}
\tauasym = 1 - \exp(-\Sigma/2).
\label{eq:tau_sigma}
\end{equation}
This single equation, applied across scales, gives:

\begin{table}[htbp]
\caption{\label{tab:time_arrow}The arrow of time at different scales.}
\begin{ruledtabular}
\begin{tabular}{lccc}
\textbf{Scale} & $\Sigma$ & $\tauasym$ & \textbf{Physical meaning} \\
\colrule
Closed system & $0$ & $0$ & No arrow of time \\
Open (lab) & $> 0$ & $> 0$ & Arrow exists \\
Near BH & $r_s/r$ & $\to 1$ & Strong arrow \\
Galaxy edge & $+ \alpha\ln r$ & $> 0$ & Arrow from RG \\
Hubble horizon & $\sim 1$ & $\sim 0.39$ & Cosmic arrow \\
Singularity & $\to \infty$ & $= 1$ & Maximal arrow \\
\end{tabular}
\end{ruledtabular}
\end{table}

\newsyn\ \emph{Three key insights.}---
\begin{enumerate}
\item \emph{$\tauasym = 0$ (closed system): time does not
  exist.}---When $\Sigma = 0$, perfect Petz recovery is
  possible: the ``future'' state can be perfectly retrodicted
  from the ``past'' state, and vice versa. Past and future are
  operationally indistinguishable. This is not a philosophical
  statement but an operational one: no measurement can
  determine the temporal ordering of events.

\item \emph{$\tauasym > 0$ (open system): time's arrow
  exists.}---When $\Sigma > 0$, retrodiction is imperfect:
  there is a preferred direction. The degree of preference is
  continuously quantified by $\tauasym \in (0, 1)$.

\item \emph{$\tauasym \to 1$ (macroscopic/singularity):
  classical world.}---When $\tauasym \approx 1$, the arrow of
  time is absolute. Collapse is complete. Classical physics
  emerges. This is the everyday world.
\end{enumerate}

\newsyn\ \emph{The signed version.}---The full picture requires
\begin{equation}
\tauasym_{\mathrm{signed}} = \ln\!\left(\frac{P_{\mathrm{fwd}}}{P_{\mathrm{rev}}}\right),
\label{eq:tau_signed}
\end{equation}
which allows three cases:
\begin{itemize}
\item $\tauasym_{\mathrm{signed}} > 0$: forward arrow (normal);
\item $\tauasym_{\mathrm{signed}} = 0$: no arrow (closed/unitary);
\item $\tauasym_{\mathrm{signed}} < 0$: reverse arrow
  (non-Markovian backflow, spin echo, Wang et~al.\ 2002
  colloidal particle experiments~\cite{Wang2002}).
\end{itemize}
The second law ($\langle \Sigma \rangle \geq 0$) is a
\emph{statistical} law: $P(\Sigma < 0)/P(\Sigma > 0) =
\exp(-|\Sigma|)$ (Crooks). Time reversal is not impossible,
merely exponentially unlikely for macroscopic systems.

%=======================================================================
\section{Cosmological Extension}
\label{sec:cosmology}
%=======================================================================

\subsection*{Content outline}

\spec\ \emph{FRW from $\Sigma$.}---In the Friedmann--Robertson--Walker
(FRW) universe, there is no timelike Killing vector. The natural
replacement is the Kodama vector $K^a$~\cite{Kodama1980}, which:
\begin{itemize}
\item Is divergence-free in any spherically symmetric spacetime;
\item Gives a conserved charge = Misner--Sharp mass;
\item Defines surface gravity at the apparent horizon.
\end{itemize}

The cosmological $\tauasym$ can be defined via~\cite{Paper3cosmological}:
\begin{equation}
F_K = \frac{|K|}{|K_{\mathrm{ref}}|},
\label{eq:F_kodama}
\end{equation}
giving $\tauasym = 0$ at the observer ($r = 0$) and
$\tauasym = 1$ at the apparent horizon ($|K|^2 = 0$).

\spec\ \emph{Friedmann equations from QRE.}---The first law of
thermodynamics at the FRW apparent horizon,
$dE = TdS + WdV$, projects along the Kodama vector to give the
second Friedmann equation~\cite{Hayward1998,CaiKim2005}.
Combined with Dorau--Much and Bianconi, this suggests:
\begin{equation}
H^2 + \frac{k}{a^2} = \frac{8\pi G}{3}\,\rho
\quad \Leftarrow \quad
\delta\, D(\rho_{\mathrm{st}} \| \rho_{\mathrm{m}}) = 0.
\label{eq:friedmann_from_qre}
\end{equation}

\toderive{Derive the Friedmann equations from
$\Sigma = D(\rho_{\mathrm{spacetime}} \| \rho_{\mathrm{matter}})$
using Bianconi's GfE framework with FRW boundary conditions.
Bianconi~\cite{BianconiCosmology2025} has shown FRW solutions
exist; the task is to reproduce the standard Friedmann equations
as a limit.}

\spec\ \emph{Dark energy from finite cosmic information.}---
Padmanabhan's CosMIn principle~\cite{Padmanabhan2017}: the total
cosmic information $I_{\mathrm{CosMIn}}$ transferred from
quantum to classical must be \emph{finite}. This requires
late-time accelerated expansion ($\Lambda > 0$). In the
$\tauasym$ language:
\begin{equation}
I_{\mathrm{CosMIn}} = \int_0^\infty \tauasym(t)\, dt < \infty
\quad \Rightarrow \quad \Lambda > 0.
\label{eq:cosmin_tau}
\end{equation}

\towrite{Paragraph on the interpretation: $\Lambda > 0$ is not
a ``mystery'' but a \emph{necessity} for finite cosmic information
processing. The universe must eventually stop producing
information (entropy) at an unbounded rate. Accelerated
expansion achieves this by diluting matter and radiation.}

%=======================================================================
\section{The Analogy Made Precise}
\label{sec:analogy}
%=======================================================================

\subsection*{Content outline}

\newsyn\ \emph{One equation, three solutions.}---We now make
the $F = ma$ analogy precise:

\begin{table*}[htbp]
\caption{\label{tab:analogy}The $F = ma$ analogy for
$\Sigma = D(\rho_{\mathrm{st}} \| \rho_{\mathrm{m}})$.}
\begin{ruledtabular}
\begin{tabular}{p{3cm}p{3.5cm}p{3.5cm}p{3.5cm}}
& \textbf{$F = ma$} & \textbf{$\Sigma = D(\rho_{\mathrm{st}} \| \rho_{\mathrm{m}})$} & \textbf{Status} \\
\colrule
Equation & $\mathbf{F} = m\mathbf{a}$ & $\Sigma = D(\rho_{\mathrm{st}} \| \rho_{\mathrm{m}})$ & \spec \\[4pt]
Solution 1 & Kepler orbits (gravity, $r \gg 0$) & Exponential metric (strong field, $r \sim r_s$) & \newsyn \\[4pt]
Solution 2 & SHM (Hooke, spring) & Flat rotation curves (weak field, $r \sim$ kpc) & \newsyn \\[4pt]
Solution 3 & Navier--Stokes (continuum) & Standard QI ($G = 0$, any CPTP) & \known \\[4pt]
Solution 4 & --- & FRW cosmology ($r \sim L_H$) & \spec \\[4pt]
Matching & Virial theorem & $\Sigma_{\mathrm{strong}} \to \Sigma_{\mathrm{weak}}$ at $r \gg r_s$ & \toderive{} \\[4pt]
Deep structure & Lagrangian mechanics & Petz recovery map as retrodiction functor & \known \\
\end{tabular}
\end{ruledtabular}
\end{table*}

\emph{What the analogy captures.}---
\begin{itemize}
\item Just as $F = ma$ does not tell you the force law (you
  need to specify gravity, Hooke, etc.), $\Sigma = D(\rho_{\mathrm{st}}
  \| \rho_{\mathrm{m}})$ does not tell you the channel (you need
  to specify the boundary conditions).
\item Just as different solutions of $F = ma$ look completely
  different (ellipses vs.\ oscillations vs.\ turbulence), different
  solutions of the unified $\Sigma$ look completely different
  (black holes vs.\ rotation curves vs.\ quantum channels).
\item Just as $F = ma$ has a deeper formulation (Lagrangian
  mechanics, Hamilton's principle), the unified $\Sigma$ has
  a deeper formulation via the Petz recovery map as the canonical
  retrodiction functor~\cite{Parzygnat2023}.
\end{itemize}

%=======================================================================
\section{Open Problems}
\label{sec:open}
%=======================================================================

\subsection*{Content outline}

\towrite{Each open problem should be 2--3 paragraphs: statement,
current status, possible approach.}

\emph{Problem 1: First-principles derivation.}---Derive
$\Sigma = D(\rho_{\mathrm{spacetime}} \| \rho_{\mathrm{matter}})$
from more primitive axioms. Currently, Eq.~\eqref{eq:unified} is
a \emph{conjecture} motivated by Bianconi's GfE~\cite{Bianconi2025},
Dorau--Much's QRE $\to$ Einstein equations~\cite{DorauMuch2025},
and the consistency of Papers~1--3. A derivation from axioms
(e.g., von~Neumann algebra axioms, or the axioms of algebraic
QFT) would elevate it from conjecture to theorem.

\spec\ Current best candidates: (i)~Bianconi's entropic action
principle~\cite{Bianconi2025}, though its vacuum solutions give
Schwarzschild $+ O(r^{-4})$ corrections rather than the exponential
metric~\cite{Thattarampilly2026}, so it cannot serve as the
mathematical realization in its current form;
(ii)~Dorau--Much's modular
theory~\cite{DorauMuch2025} extended to include backreaction;
(iii)~Jacobson's entanglement equilibrium~\cite{Jacobson2016}
rewritten in QRE language. The three action-independent routes
(modular flow, gravitational Landauer, quantum channel) may
provide the necessary foundation without requiring a specific
modified gravity action.

\emph{Problem 2: Planck scale.}---What happens when
spacetime itself is quantum? At the Planck scale,
$\rho_{\mathrm{spacetime}}$ is not a semiclassical density
matrix but a fully quantum object. The QRE
$D(\rho_{\mathrm{spacetime}} \| \rho_{\mathrm{matter}})$ must
be replaced by an algebraic QFT version using the Araki
relative entropy~\cite{Araki1976}. This is mathematically
well defined but computationally intractable with current
methods.

\emph{Problem 3: Cosmological initial conditions.}---Why does
the universe start with $\Sigma \approx 0$ (low entropy)?
In the $\tauasym$ framework, this is the question: why does
$\tauasym_{\mathrm{initial}} \approx 0$? One possibility:
the initial state is a ``Petz-recovered'' state from the
final condition (a closed-universe scenario where boundary
conditions at both ends enforce $\Sigma = 0$). This is
highly speculative and connects to the ``past hypothesis''
of thermodynamics~\cite{Albert2000,Carroll2010}.

\emph{Problem 4: The Born rule.}---The $\tauasym$ framework
assumes the Born rule in computing fidelity $F$ and entropy
production $\Sigma$. A complete theory should derive the
Born rule, not assume it. Possible connections:
\begin{itemize}
\item Zurek's ``envariance'' derivation~\cite{Zurek2005}
  uses environment-assisted symmetry, which connects to
  $\tauasym$ via quantum Darwinism;
\item Deutsch--Wallace decision-theoretic approach relies
  on rational betting, which maps to Petz recovery
  optimality;
\item Gleason's theorem constrains probability measures on
  Hilbert space, and the Petz map is the unique ``probability
  measure on channels.''
\end{itemize}

\spec\ \emph{Problem 5: Consciousness.}---We acknowledge but do
not attempt this problem. The hard problem of consciousness
(why subjective experience exists) is sometimes connected to
the arrow of time and measurement. The $\tauasym$ framework
provides a precise, quantitative arrow of time that does not
require an observer. Whether this has implications for
consciousness is a question for philosophers, not for this paper.

\emph{Problem 6: Derive $\eta = 1$.}---The marginal anomalous
dimension $\eta = 1$ that produces flat rotation curves
(Paper~3) should follow from the properties of
$D(\rho_{\mathrm{spacetime}} \| \rho_{\mathrm{matter}})$ in
the IR. This is the most concrete open problem and the one
most likely to be resolved in the near term. The DPI
structure of QRE may constrain $\eta$ to exactly 1 via
a saturation argument.

\emph{Problem 7: Gravitational waves and dynamical spacetimes.}---
Papers~2 and~3 address static or quasi-static spacetimes.
Extending to dynamical spacetimes (gravitational wave emission,
binary mergers) requires:
\begin{itemize}
\item A time-dependent $\Sigma(t)$;
\item The Kodama vector (for spherical symmetry) or its
  generalization~\cite{DorauVerch2024};
\item Raychaudhuri decomposition of $d\Sigma/ds$ along
  geodesic congruences.
\end{itemize}

%=======================================================================
\section{Conclusion}
\label{sec:conclusion}
%=======================================================================

\subsection*{Content outline}

\towrite{Concluding section synthesizing the full four-paper arc.}

\emph{Summary.}---We have argued that a single equation,
$\Sigma = D(\rho_{\mathrm{spacetime}} \| \rho_{\mathrm{matter}})$,
unifies three scales of physics:
\begin{itemize}
\item Quantum information (Paper~1): $\Sigma = 0$
  $\Leftrightarrow$ perfect recovery, no arrow of time;
\item Strong-field gravity (Paper~2): $\Sigma = r_s/r$
  $\to$ exponential metric, information-theoretic ``soft horizon'';
\item Galactic dynamics (Paper~3): $\Sigma = r_s/r + \alpha\ln r$
  $\to$ running $G$, flat rotation curves without dark matter
  particles.
\end{itemize}

\emph{The philosophical point.}---
\begin{quote}
We have not discovered a new equation. We have proposed that an
existing mathematical structure---quantum relative entropy and
its recovery theory---when applied consistently across scales,
unifies phenomena that appear unrelated. The contribution is not
a new formula but a new viewpoint that makes existing results
cohere. This is the kind of contribution that, if correct,
simplifies rather than complicates.
\end{quote}

\towrite{Final paragraph on falsifiability: the framework makes
specific, testable predictions at each scale (Paper~2: GW echoes,
shadow sizes; Paper~3: universal $r_c$, SPARC fits, BIG-SPARC).
If these predictions fail, the framework fails. This is a
feature, not a bug.}

\begin{acknowledgments}
The author thanks the quantum information, quantum gravity,
astrophysics, and quantum foundations communities for the
foundational works on which this synthesis builds.
\end{acknowledgments}

%-----------------------------------------------------------------------
% BIBLIOGRAPHY
%-----------------------------------------------------------------------

\begin{thebibliography}{60}

\bibitem{Paper1}
S.-K. Huang,
``Universal recovery: The Petz map as retrodiction functor,
error corrector, and entropy bound,''
companion paper (2026).

\bibitem{Paper2}
S.-K. Huang,
``Information-theoretic origin of the exponential metric:
Gravitational redshift as Petz recovery fidelity,''
companion paper (2026).

\bibitem{Paper3}
S.-K. Huang,
``Flat rotation curves from quantum relative entropy:
Running Newton's constant as information processing,''
companion paper (2026).

\bibitem{Bianconi2025}
G.~Bianconi,
``Gravity from entropy,''
Phys.\ Rev.\ D \textbf{111}, 066001 (2025);
arXiv:2408.14391.

\bibitem{BianconiCosmology2025}
G.~Bianconi,
``Gravity from entropy: Cosmological solutions,''
(2025); arXiv:2510.22545.

\bibitem{DorauMuch2025}
M.~Dorau and A.~Much,
``From quantum relative entropy to the semiclassical Einstein
equations,''
Phys.\ Rev.\ Lett.\ (2025); arXiv:2510.24491.

\bibitem{CasiniHuerta2017}
H.~Casini and M.~Huerta,
``Relative entropy and the RG flow,''
J.\ High Energy Phys.\ \textbf{2017}, 044 (2017);
arXiv:1611.00016.

\bibitem{Kumar2025}
N.~Kumar,
``Marginal IR running of gravity as a natural explanation for
dark matter,''
Phys.\ Lett.\ B \textbf{871}, 140008 (2025);
arXiv:2509.05246.

\bibitem{Gubitosi2024}
G.~Gubitosi, O.~F. Piattella, and R.~M. Casarini,
``Phenomenology of renormalization group improved gravity from
the kinematics of SPARC galaxies,''
(2024); arXiv:2403.00531.

\bibitem{Parzygnat2023}
A.~J. Parzygnat and F.~Buscemi,
``Axioms for retrodiction: Achieving time-reversal symmetry with a
prior,''
Quantum \textbf{7}, 1013 (2023).

\bibitem{Jacobson1995}
T.~Jacobson,
``Thermodynamics of spacetime: The Einstein equation of state,''
Phys.\ Rev.\ Lett.\ \textbf{75}, 1260 (1995).

\bibitem{Jacobson2016}
T.~Jacobson,
``Entanglement equilibrium and the Einstein equation,''
Phys.\ Rev.\ Lett.\ \textbf{116}, 201101 (2016);
arXiv:1505.04753.

\bibitem{Araki1976}
H.~Araki,
``Relative entropy of states of von Neumann algebras,''
Publ.\ RIMS Kyoto \textbf{11}, 809 (1976).

\bibitem{Kodama1980}
H.~Kodama,
``Conserved energy flux for the spherically symmetric system and
the backreaction problem in the black hole evaporation,''
Prog.\ Theor.\ Phys.\ \textbf{63}, 1217 (1980).

\bibitem{Hayward1998}
S.~A. Hayward,
``Unified first law of black-hole dynamics and relativistic
thermodynamics,''
Class.\ Quantum Grav.\ \textbf{15}, 3147 (1998).

\bibitem{CaiKim2005}
R.-G. Cai and S.~P. Kim,
``First law of thermodynamics and Friedmann equations of
Friedmann--Robertson--Walker universe,''
J.\ High Energy Phys.\ \textbf{2005}, 050 (2005).

\bibitem{Padmanabhan2017}
T.~Padmanabhan,
``Cosmic information, the cosmological constant, and the
amplitude of primordial perturbations,''
Phys.\ Lett.\ B \textbf{773}, 81 (2017);
arXiv:1703.06144.

\bibitem{Verlinde2016}
E.~P. Verlinde,
``Emergent gravity and the dark universe,''
SciPost Phys.\ \textbf{2}, 016 (2017); arXiv:1611.02269.

\bibitem{Smolin2017}
L.~Smolin,
``MOND as a regime of quantum gravity,''
Phys.\ Rev.\ D \textbf{96}, 083523 (2017);
arXiv:1704.00780.

\bibitem{Wang2002}
G.~M. Wang, E.~M. Sevick, E.~Mittag, D.~J. Searles,
and D.~J. Evans,
``Experimental demonstration of violations of the second law
of thermodynamics for small systems and short time scales,''
Phys.\ Rev.\ Lett.\ \textbf{89}, 050601 (2002).

\bibitem{Zurek2005}
W.~H. Zurek,
``Probabilities from entanglement, Born's rule
$p_k = |\psi_k|^2$ from envariance,''
Phys.\ Rev.\ A \textbf{71}, 052105 (2005).

\bibitem{Albert2000}
D.~Z. Albert,
\emph{Time and Chance}
(Harvard University Press, 2000).

\bibitem{Carroll2010}
S.~Carroll,
\emph{From Eternity to Here: The Quest for the Ultimate Theory
of Time}
(Dutton, 2010).

\bibitem{ReuterWeyer2004}
M.~Reuter and H.~Weyer,
``Running Newton constant, improved gravitational actions, and
galaxy rotation curves,''
Phys.\ Rev.\ D \textbf{70}, 124028 (2004).

\bibitem{SPARC2016}
F.~Lelli, S.~S. McGaugh, and J.~M. Schombert,
``SPARC: Mass models for 175 disk galaxies with Spitzer
photometry and accurate rotation curves,''
Astron.\ J.\ \textbf{152}, 157 (2016); arXiv:1606.09251.

\bibitem{McGaugh2016}
S.~S. McGaugh, F.~Lelli, and J.~M. Schombert,
``Radial acceleration relation in rotationally supported
galaxies,''
Phys.\ Rev.\ Lett.\ \textbf{117}, 201101 (2016);
arXiv:1609.05917.

\bibitem{GhariHaghi2026}
A.~Ghari and H.~Haghi,
``Testing emergent gravity in dwarf spheroidal galaxies,''
(2026); arXiv:2601.01715.

\bibitem{BrouwerKiDS2021}
M.~M. Brouwer \emph{et~al.},
``The weak lensing radial acceleration relation,''
Astron.\ Astrophys.\ \textbf{650}, A113 (2021);
arXiv:2106.11677.

\bibitem{Petz1986}
D.~Petz,
``Sufficient subalgebras and the relative entropy of states of a
von Neumann algebra,''
Commun.\ Math.\ Phys.\ \textbf{105}, 123 (1986).

\bibitem{BassoMazieroCeleri2025}
M.~L. Basso, J.~Maziero, and L.~C. C\'{e}leri,
``Quantum detailed fluctuation theorem in curved spacetime,''
Phys.\ Rev.\ Lett.\ \textbf{134}, 050406 (2025);
arXiv:2405.03902.

\bibitem{DorauVerch2024}
M.~Dorau and R.~Verch,
``Kodama-like vectors in axisymmetric spacetimes,''
(2024); arXiv:2402.18993.

\bibitem{Clowe2006}
D.~Clowe \emph{et~al.},
``A direct empirical proof of the existence of dark matter,''
Astrophys.\ J.\ Lett.\ \textbf{648}, L109 (2006).

\bibitem{Donoghue2020}
J.~F. Donoghue,
``A critique of the asymptotic safety program,''
Front.\ Phys.\ \textbf{8}, 56 (2020).

\bibitem{Paper3cosmological}
S.-K. Huang,
research notes on cosmological $\tauasym$ (2026).

\bibitem{Thattarampilly2026}
B.~Thattarampilly, Y.~Zheng, and K.~K.~Kakkat,
``Spherically symmetric black hole in gravity from entropy,''
(2026); arXiv:2602.13694.

\bibitem{Herrera2020}
L.~Herrera,
``The Gibbs paradox, the Landauer principle and the irreversibility
associated with tilted observers,''
Entropy \textbf{22}, 340 (2020).

\bibitem{ChenSongScarani2025}
Z.~Chen, Y.~Song, and V.~Scarani,
``Petz recovery from subsystems in bosonic quantum error correction,''
(2025); arXiv:2511.05941.

\end{thebibliography}

\end{document}
